%==============================================================================
%  07-reinvention-roadmap.tex — Section 6: Strategic Reinvention Roadmap
%==============================================================================
\strategyPart{19 \textbar\ Transformation Roadmap}{The Strategic Reinvention Roadmap}

Transformation is a change-management problem, not a technology problem. The
most widely validated change framework — Kotter's eight-step model of leading
change [Kotter Inc.] — maps directly onto AI-enabled reinvention. This section
compresses the eight steps into five logical phases that a CEO can sequence
without depending on calendar time.

\section{Kotter's eight steps}

\begin{itemize}
  \item \textbf{1 · Create a sense of urgency} — a compelling, quantified case for why AI reinvention cannot wait.
  \item \textbf{2 · Build a guiding coalition} — a cross-functional team with the authority to act.
  \item \textbf{3 · Form a strategic vision} — a clear picture of the reinvented organization.
  \item \textbf{4 · Enlist a volunteer army} — mobilize the broader organization behind the vision.
  \item \textbf{5 · Enable action by removing barriers} — dismantle silos, data walls and legacy rules.
  \item \textbf{6 · Generate short-term wins} — visible, measurable successes early.
  \item \textbf{7 · Sustain acceleration} — build on credibility to scale across the enterprise.
  \item \textbf{8 · Institute change} — anchor new behaviors in culture, systems and incentives.
\end{itemize}

\section{The five phases of AI reinvention}

\begin{center}
  \renewcommand{\arraystretch}{1.5}
  \rowcolors{2}{inSnow}{white}
  \fitwidth{\begin{tabular}{>{\raggedright\arraybackslash}p{2.2cm}>{\raggedright\arraybackslash}p{4.3cm}>{\raggedright\arraybackslash}p{5.1cm}}
    \rowcolor{inNavy}
    \hcell{Phase} & \hcell{Core work} & \hcell{Kotter steps applied} \\
    \textbf{1 · Diagnosis} & Map the highest-value organizational problems and audit the three stocks. & 1 -- Urgency \\
    \textbf{2 · Cultural foundation} & Literacy, trust and a guiding coalition across the C-suite and functions. & 2 -- Coalition; 3 -- Vision \\
    \textbf{3 · Proof of value} & Targeted Assist/Automate deployments with hard KPIs and visible wins. & 4 -- Volunteer army; 5 -- Remove barriers; 6 -- Short-term wins \\
    \textbf{4 · Scaling \& redesign} & Workflow redesign and Agentic/Optimization deployments at enterprise scale. & 7 -- Sustain acceleration \\
    \textbf{5 · Reinvention} & Business-model reinvention anchored in culture and incentives. & 8 -- Institute change \\
  \end{tabular}}
\end{center}

\section{The seven-phase roadmap from traditional to AI-based}

\begin{center}
  \renewcommand{\arraystretch}{1.45}
  \rowcolors{2}{inSnow}{white}
  \fitwidth{\begin{tabular}{>{\raggedright\arraybackslash}p{2.6cm}>{\raggedright\arraybackslash}p{4.6cm}>{\raggedright\arraybackslash}p{4.4cm}}
    \rowcolor{inNavy}
    \hcell{Phase} & \hcell{Key activities} & \hcell{Main risk} \\
    \textbf{1 · Diagnosis} & Maturity and readiness assessment & Incomplete assessment \\
    \textbf{2 · Strategic alignment} & AI vision, target operating model, priorities & Technology-led strategy \\
    \textbf{3 · Controlled experimentation} & Limited, scored use cases & Pilot without scale \\
    \textbf{4 · Workflow redesign} & Process and role redesign & Employee resistance \\
    \textbf{5 · Platform \& capability} & Data, model gateway, evaluation, monitoring & Inconsistent standards \\
    \textbf{6 · Scaling} & Reusable patterns, AI product teams, funding & Governance gaps \\
    \textbf{7 · Operating-model reinvention} & Structures, decision rights, business model & Governance complexity \\
  \end{tabular}}
\end{center}

\vspace{0.6em}
\begin{center}
  \fitwidth{\begin{tikzpicture}
    \node[fill=inNavy,   text=white, font=\bfseries\footnotesize, minimum height=0.8cm, text width=2.1cm, align=center, inner sep=2pt] (p1) {Phase 1\\Diagnosis};
    \node[fill=inNavy,   text=white, font=\bfseries\footnotesize, minimum height=0.8cm, text width=2.1cm, align=center, inner sep=2pt, right=2pt of p1] (p2) {Phase 2\\Alignment};
    \node[fill=inHover,  text=white, font=\bfseries\footnotesize, minimum height=0.8cm, text width=2.1cm, align=center, inner sep=2pt, right=2pt of p2] (p3) {Phase 3\\Experiments};
    \node[fill=inHover,  text=white, font=\bfseries\footnotesize, minimum height=0.8cm, text width=2.1cm, align=center, inner sep=2pt, right=2pt of p3] (p4) {Phase 4\\Redesign};
    \node[fill=inBright, text=white, font=\bfseries\footnotesize, minimum height=0.8cm, text width=2.1cm, align=center, inner sep=2pt, right=2pt of p4] (p5) {Phase 5\\Platform};
    \node[fill=inBright, text=white, font=\bfseries\footnotesize, minimum height=0.8cm, text width=2.1cm, align=center, inner sep=2pt, right=2pt of p5] (p6) {Phase 6\\Scaling};
    \node[fill=inLight,  text=white, font=\bfseries\footnotesize, minimum height=0.8cm, text width=2.1cm, align=center, inner sep=2pt, right=2pt of p6] (p7) {Phase 7\\Reinvention};
  \end{tikzpicture}}
\end{center}

\section{Sequencing principles}

\begin{itemize}
  \item \textbf{Problems first, proofs second.} Phase 1 is a diagnostic, not a
        procurement exercise; it identifies the problems before any tool is
        chosen.
  \item \textbf{Every phase ends with a measurable checkpoint.} Short-term wins
        are the fuel of Kotter's model and the antidote to the 30\% proof-of-
        concept abandonment rate [Gartner 2023].
  \item \textbf{Culture is not a phase to ``get through.''} The second phase
        builds the literacy and trust that every later phase depends on
        [HBR 2018].
\end{itemize}
